% Strategic Management - Week 14 Cornell Final Notes (Spring 2026, Wed/Thu)
% Compile: pdflatex strategic-management-final-cornell-master-notes-wed-thu-spring2026.tex
\documentclass[10pt,a4paper]{article}
\usepackage[utf8]{inputenc}
\usepackage[T1]{fontenc}
\usepackage[margin=1.2cm]{geometry}
\usepackage{xcolor}
\usepackage{longtable}
\usepackage{enumitem}
\setlist[itemize]{leftmargin=1.2em,itemsep=1pt,topsep=2pt}
\definecolor{sblue}{RGB}{24,52,110}
\definecolor{soft}{RGB}{245,248,252}
\setlength{\parindent}{0pt}
\setlength{\parskip}{0.35em}

\newcommand{\cornell}[3]{
\subsection*{#1}
\begin{longtable}{|p{0.22\linewidth}|p{0.74\linewidth}|}
\hline
\rowcolor{soft}\textbf{Cue / Questions} & \textbf{Notes} \\
\hline
#2 & #3 \\
\hline
\end{longtable}
}

\begin{document}
\begin{center}
{\Large\bfseries\color{sblue} Strategic Management - Week 14 Final Cornell Master Notes}\\
{\small Spring 2026 (typical class days: Wed/Thu) | Consolidated from Units 0-6 notes, class discussions, case work, and midterm prep}
\end{center}

\section*{Performance and Focus Rules}
\begin{itemize}
  \item \textbf{No separate numeric target provided here in your latest note}, so maintain high output quality and transfer this method to the high-risk courses.
  \item \textbf{Related global rules:} WIP fail 3.49--0; UCM fail 4.99--0; suspenso = 4.
  \item \textbf{Use this course as model discipline:} structured frameworks, mechanism logic, and evidence-led case answers.
\end{itemize}

\section*{Cornell Notes by Unit Integration}

\cornell{1) Strategy Foundations (Units 0-1)}
{What is strategy?\\
What makes decision strategic?\\
How do deliberate/emergent logics combine?}
{
\textbf{Strategy core:} long-term direction under uncertainty with scarce resources and trade-offs.\\
\textbf{Strategic decision traits:} high irreversibility, system-level consequences, and opportunity-cost implications.\\
\textbf{Process duality:} rational planning and emergent adaptation coexist in real organizations.\\
\textbf{Responsibility map:} board, CEO, TMT, and supporting analysis roles interact.\\
\textbf{Exam line:} ``Strong strategy aligns goals, environment understanding, internal resources, and implementation quality.''
}

\cornell{2) Objectives, Stakeholders, and Governance (Unit 2)}
{Mission vs vision?\\
How are objectives measured?\\
Why governance matters?}
{
\textbf{Mission/vision distinction:} present identity vs long-term aspiration.\\
\textbf{Objective design:} attribute + metric + target + timeframe.\\
\textbf{Performance lenses:} accounting, market, and social-environmental measures.\\
\textbf{Governance logic:} agency tensions require incentives, oversight, and accountability.\\
\textbf{Critical framing:} shareholder and stakeholder perspectives imply different strategic constraints.
}

\cornell{3) External Analysis (Unit 3)}
{PESTEL role?\\
How do Five Forces determine attractiveness?\\
What is industry vs firm effect?}
{
\textbf{PESTEL function:} identify macro opportunities/threats outside firm control.\\
\textbf{Five Forces purpose:} evaluate structural profit potential and pressure points.\\
\textbf{Interpretive rule:} stronger forces generally reduce average industry returns.\\
\textbf{Nuance:} firm effect can still produce superior outcomes in difficult industries.\\
\textbf{Applied insight:} clusters and institutional context can alter practical competitive dynamics.
}

\cornell{4) Internal Analysis and RBV (Unit 4)}
{Where is value created?\\
Which resources are strategic?\\
What sustains advantage?}
{
\textbf{Value chain lens:} locate primary/support activities and coordination links.\\
\textbf{RBV logic:} heterogenous resources and capabilities drive persistent performance differences.\\
\textbf{VRIO test:} valuable, rare, hard to imitate, and organized for capture.\\
\textbf{Sustainability conditions:} durability, low substitutability, low transferability, complementarity.\\
\textbf{Case-use rule:} map capabilities to mechanism and then to observable performance.
}

\cornell{5) Competitive and Corporate Strategy (Units 5-6)}
{Cost vs differentiation?\\
How does strategy clock extend Porter?\\
Where to compete at corporate level?}
{
\textbf{Unit 5 core:} competitive advantage requires coherent positioning and renewal.\\
\textbf{Cost leadership/differentiation:} both need contextual fit and risk management.\\
\textbf{Strategy clock value:} richer price-value positioning for hybrid and nuanced cases.\\
\textbf{Unit 6 core:} corporate strategy answers scope question (specialization, diversification, integration, restructuring).\\
\textbf{Integration sentence:} ``Business strategy explains how to compete; corporate strategy explains where to compete; durable performance needs both aligned.''
}

\cornell{6) Case Method (Starbucks, Wal-Mart Style Prompts)}
{How to structure case answers?\\
How to avoid descriptive-only responses?\\
What earns top marks?}
{
\textbf{Recommended structure:} issue -> framework -> evidence -> strategic option -> risk and mitigation.\\
\textbf{Evidence discipline:} use case facts to validate mechanism claims.\\
\textbf{Renewal lens:} advantage persistence must be demonstrated, not assumed.\\
\textbf{Common weakness:} naming frameworks without causal reasoning.\\
\textbf{High-score trait:} explicit trade-offs and coherent recommendation under constraints.
}

\cornell{7) Lecture Introduction: Structure, Grading, Philosophy}
{How is the course assessed?\\
How does the professor mix theory with real cases?\\
Why is strategy a ''social science''?}
{
\textbf{Grading:} 50\% Exam (theory + short questions) and 50\% Business Cases (group + individual).\\
\textbf{Methodology:} theory lectures followed by practical ``Illustrations'' and case studies (e.g., Madonna, Inditex/Zara).\\
\textbf{Strategy as social science:} not pure math---intuition, risk, and human complexity shape outcomes.\\
\textbf{Exam implication:} connect academic concepts to what managers actually do under constraints and uncertainty.
}

\cornell{8) What Strategy Means (Etymology and Core Competition)}
{What is the word ``strategy'' telling you?\\
What is the battlefield idea?}
{
\textbf{Etymology:} from Greek \textit{strategos} (general commanding an army).\\
\textbf{Core concept:} strategy is competing for scarce resources---customers, talent, and capital---in a market ``battlefield''.\\
\textbf{Decision horizon logic:} the sequence of choices becomes harder to reverse as time passes (strategic decisions create path dependence).
}

\cornell{9) Levels of Decision-Making (Strategic / Tactical / Operational)}
{What changes across levels?\\
How do the corporate/business/functional layers relate?}
{
\textbf{Strategic:} long-term, high-impact choices (e.g., degree choice, buying a house, children).\\
\textbf{Tactical:} medium-term coordination (align departments).\\
\textbf{Operational:} day-to-day execution.\\
\textbf{Hierarchy across firm layers:} corporate level (\textit{where we want to be}) -> business level (\textit{how we compete}) -> functional level (\textit{how we support the business}).
}

\cornell{10) Why Strategies Fail + The ``Know Thyself'' Rule}
{Why do plans collapse even with good frameworks?\\
What must managers analyze first?}
{
\textbf{Most failures come from:} \textbf{Irrationality} (family-owned biases, unions, politics) or \textbf{External volatility} (perfect forecasting is impossible beyond ~2 years).\\
\textbf{Know thyself rule:} internal analysis before facing competitors; map \textbf{strengths and weaknesses} like a general understanding the limits of their army.\\
\textbf{Exam move:} explain the failure mechanism---not only the conclusion.
}

\cornell{11) Sustainability & Ethics (Greenwashing to Social Impact)}
{How does sustainability enter strategy?\\
What tension exists between profit logic and social logic?}
{
\textbf{Greenwashing:} deceptive marketing that claims environmental friendliness without substantive change.\\
\textbf{Shift in framing:} from pure profit-seeking to considering societal impact (stakeholder reality).\\
\textbf{Real-world conflict:} \textbf{Economic vs. Social} priorities---Milton Friedman-style shareholder profit logic vs the professor's view that family/worker protection/geopolitics often override simplified rational economics.\\
\textbf{Strategic takeaway:} success requires both market knowledge and deep awareness of internal limits and constraints.
}

\cornell{12) SWOT Foundation (PESTEL + the Strategic Mountain)}
{What does SWOT actually analyze?\\
How do opportunities/threats connect to PESTEL?}
{
\textbf{Internal analysis:} Strengths/Weaknesses (resources and capabilities).\\
\textbf{External analysis:} Opportunities/Threats (use \textbf{PESTEL}: Political, Economic, Social, Technological, Environmental, Legal).\\
\textbf{Purpose:} guide strategic choices by linking environment pressures to what the firm can realistically do.\\
\textbf{Case-use pattern:} for each SWOT item, state the mechanism and the implication for strategic options.
}

\cornell{13) Vocabulary Word Bank (Exam Anchors)}
{What definitions should be ready for short questions?\\
Which examples are recurring in cases?}
{
\textbf{Strategic Management:} continuous planning, monitoring, and assessment of what is necessary to meet an organization's goals.\\
\textbf{Etymology:} origin of words (e.g., \textit{strategy} from Greek \textit{strategia/strategos}).\\
\textbf{Forecasting:} predicting future business trends; professor warns it is rarely accurate for long horizons (often beyond \textasciitilde 2 years).\\
\textbf{Greenwashing:} deceptive marketing that signals environmental friendliness without real substance.\\
\textbf{FDI (Foreign Direct Investment):} controlling ownership in a business in another country by an entity based abroad.\\
\textbf{Inditex (Zara):} Spanish multinational clothing company used as a case lens for supply-chain execution and market dominance.
}

\cornell{14) Field Evolution + The Four Pillars (Grant Model)}
{How did the field develop over time?\\
What makes a strategy actually work?}
{
\textbf{Field evolution:} 1960s \textit{Business Policy} -> 1970s \textit{Strategic Planning} -> current \textit{Strategic Management} (adaptation + execution emphasis).\\
\textbf{Core definition:} long-term objectives + a ''course of action'' that achieves them despite scarce resources and competition.\\
\textbf{Four pillars (Grant-style balance):}\\
1) simple, consistent, long-term goals (clarity on what we want to be)\\
2) understanding the competitive environment (''rules of the game'', rivals)\\
3) objective appraisal of resources (strengths + weaknesses; SWOT discipline)\\
4) effective implementation (best plan fails if poorly executed or without organizational commitment).
}

\cornell{15) Levels of Strategy + The Human Element}
{Where does ''strategy'' live inside an organization?\\
Who actually shapes strategic outcomes?}
{
\textbf{Corporate level:} where to compete (scope of industries/geographies).\\
\textbf{Business level:} how to compete (positioning vs rivals; competitive advantage logic).\\
\textbf{Functional level:} how to execute (HR, marketing, finance, operations).\\
\textbf{Human element:} strategies are built by CEOs/managers; outcomes reflect bias, age/education, incentives, and ego (not ''legal fictions'').
}

\cornell{16) Agency Theory (Owners vs Managers)}
{What conflict changes strategic incentives?\\
Why can ''profit-maximizing'' goals fail?}
{
\textbf{Principal-agent tension:} shareholders (owners) want profit/max dividends; managers can pursue different utility.\\
\textbf{Manager utility (typical):} status, power, and higher remuneration.\\
\textbf{Strategic implication:} growth via mergers/acquisitions can be pursued even when it does not create value for shareholders.\\
\textbf{Exam phrasing:} explain the incentive mismatch and then connect it to the observed strategic choice (e.g., ''empire building'').
}

\cornell{17) Forecasting Limits + Prospection Scenarios}
{Why do long-term ''predictions'' break?\\
How should managers handle uncertainty?}
{
\textbf{Key rule:} accurate forecasting is hard beyond about \textasciitilde 2 years because of ''black swan'' events (geopolitics, climate, health shocks).\\
\textbf{Manager response:} prospection via scenarios; plan for plausible futures rather than one forecast.
}

\cornell{18) Strategic Decision-Context Statistics (Family + M\&A)}
{How common are ''real'' organizational structures?\\
What do stats imply for strategy design?}
{
\textbf{Family-owned firms:} roughly 80--90\% of businesses worldwide; strategic choices often filtered through family values and ''non-economic'' constraints.\\
\textbf{M\&A prevalence and risk:} many mergers fail to create shareholder value; acquirers often pay a 20--50\% premium to win bids.\\
\textbf{Mechanism:} premium and deal logic can be driven by CEO desire for a larger empire, not by value-creation math.
}

\cornell{19) Vocabulary Word Bank (PESTEL + Integration Terms)}
{What terms should be instant on exams?\\
What do they mean for strategy options?}
{
\textbf{PESTEL analysis:} Political, Economic, Social, Technological, Environmental, and Legal macro forces shaping the environment.\\
\textbf{Synergy:} combined value/performance of two firms greater than the sum of separate parts.\\
\textbf{Vertical integration:} expansion along steps of the same production path (e.g., manufacturer buying a supplier).\\
\textbf{Diversification:} corporate strategy to enter new markets/industries outside current business scope.
}

\cornell{20) Realist Strategy: Mission, Vision, Objectives}
{What is mission vs vision?\\
What are strategic objectives?\\
How does strategy bridge the present to a 10--20 year future?}
{
\textbf{Realist framing:} firms are social realities made of people with conflicting motivations (stakeholders vs shareholders), not just abstract legal entities.\\
\textbf{Mission:} ''Who am I now?'' (current identity and purpose).\\
\textbf{Vision:} ''Where do we want to be in 10--20 years?'' (aspirational future).\\
\textbf{Strategic gap:} distance between mission (present reality) and vision (future reality).\\
\textbf{Strategic objectives:} 6-month to 1.5-year steps that bridge the gap with measurable targets.\\
\textbf{Bridge logic:} strategy connects today's constraints and conflicts to future direction through objectives and execution.
}

\cornell{21) Success Metrics: Accounting Profit vs ROA/ROE}
{How do you measure success?\\
Why are absolute numbers misleading?\\
How do opportunity cost and agency conflicts enter?}
{
\textbf{Accounting profit:} revenue minus expenses; can be manipulated and may not reflect value creation.\\
\textbf{Profitability ratios (ROA/ROE):} relative measures (e.g., 10\% return) that enable comparison across different sizes and industries.\\
\textbf{Shareholder value:} success means creating value for owners relative to the \textbf{opportunity cost}.\\
\textbf{8\% satisfaction rule:} an 8\% return is only acceptable if the rival/best alternative is lower (e.g., 7\%); if the rival returns 12\%, you destroy value strategically.\\
\textbf{Agency layer:} owners seek profit/value; managers may prefer power/status/growth; workers may pursue better conditions and pay.
}

\cornell{22) Strategic Planning Cycle (Mission -> Objectives -> Vision)}
{What is the planning logic over time?\\
How does a person/company move from today to the future?}
{
\textbf{Mission:} defines core business and ''current purpose''.\\
\textbf{Objectives:} specific, shorter-term targets (measurable steps).\\
\textbf{Vision:} ultimate long-term aspiration.\\
\textbf{Implementation link:} objectives translate vision into actions that can be monitored, adjusted, and governed.
}

\cornell{23) Net Profit vs ROA Example (Efficiency Insight)}
{Why can Company A have lower net profit but be more successful?\\
What does ROA/ROE correct?}
{
\textbf{Illustration idea:} compare absolute net profit vs efficiency/per-unit performance.\\
\textbf{Company A:} net profit 1.0 (million) with 100 employees.\\
\textbf{Company B:} net profit 1.5 (million) with 200 employees.\\
\textbf{Reading:} net profit makes B look better, but relative profitability indicates A is more efficient (higher profit per ''unit of resources'').\\
\textbf{Exam move:} argue with relative returns (ROA/ROE) and connect back to value creation and opportunity cost.
}

\cornell{24) Vocabulary Word Bank (Stakeholders, Governance, EVA, Market Value)}
{What terms define real strategic discussions?\\
How do they show up in exam answers?}
{
\textbf{Stakeholder:} any individual/group that can affect or be affected by the firm (employees, customers, suppliers, unions).\\
\textbf{Shareholder:} legal owner; primarily concerned with return on investment (ROI) and value creation.\\
\textbf{Opportunity cost:} the forgone benefit from the next best alternative (e.g., investing elsewhere).\\
\textbf{Corporate governance:} rules and practices that direct and control the firm so managers act in owners' interests.\\
\textbf{Economic Value Added (EVA):} residual wealth measure using operating profit minus the cost of capital.\\
\textbf{Market capitalization:} total stock value = price per share x total number of shares.
}

\cornell{25) Shareholder Value, Opportunity Cost, and Satisfaction}
{How is value creation measured?\\
When is return still a strategic failure?}
{
\textbf{Shareholder value (change view):} success tied to change in market value: (current share price minus previous share price) plus dividends.\\
\textbf{Opportunity cost:} value is created only if the firm's return exceeds the return on a \textbf{similar-risk} alternative investment.\\
\textbf{Satisfaction (relative):} Case A: firm 8\%, competitor 7\% $\rightarrow$ satisfied. Case B: firm 8\%, competitor 12\% $\rightarrow$ unsatisfied (value destroyed vs the market benchmark).
}

\cornell{26) Beta (\(\beta\)) and Systematic Risk}
{What does beta measure?\\
How do you read \(\beta > 1\)?}
{
\textbf{Beta (\(\beta\)):} sensitivity of a stock's returns to the market (systematic risk proxy).\\
\textbf{Interpretation:} \(\beta = 1.2\) means roughly 20\% more volatile than the market average (moves amplify vs the index).\\
\textbf{Strategy link:} required return and ''fair'' comparisons depend on risk; profit without risk adjustment misleads owners.
}

\cornell{27) Agency Conflict and Corporate Governance (Carrots \& Sticks)}
{What do owners vs managers want?\\
How do boards and incentives align interests?}
{
\textbf{Agency conflict:} owners want maximum profit/value; managers want status, power, and high pay---expansive M\&A can inflate prestige even when it destroys shareholder value.\\
\textbf{Corporate governance:} mechanisms to constrain managerial discretion.\\
\textbf{Carrots (incentives):} bonuses, \textbf{stock options} (right to buy shares later at a preset ''strike'' price to align with share price gains).\\
\textbf{Sticks (monitoring):} board of directors oversight, reporting, and accountability.\\
\textbf{Exam line:} link incentive design to the agency problem, not only to ''fair pay.''
}

\cornell{28) Stakeholder Dynamics, Externalities, and the ''Social Reality'' Pivot}
{Who pulls the firm in different directions?\\
What costs/benefits sit outside the P\&L?}
{
\textbf{Utility map (who wants what):} owners/shareholders---profit, dividends, value creation; managers/CEOs---status, power, prestige, remuneration; workers/unions---wages, security, conditions; customers---quality at low price.\\
\textbf{Externalities:} \textbf{negative}---pollution/social costs not fully paid by the firm; \textbf{positive}---jobs, local development, tax base.\\
\textbf{Stakeholder pivot:} the firm as negotiated outcome among owners, managers, workers, and society---not a single ''pure profit'' machine.
}

\cornell{29) CSR, Sustainability, and Greenwashing (DNA vs Optics)}
{Why is CSR often criticized?\\
What counts as real sustainability?}
{
\textbf{CSR critique:} often a \textbf{voluntary} spend layered on core business rather than a structural change.\\
\textbf{Greenwashing:} small ''eco'' projects for image while the core model stays unchanged.\\
\textbf{Shift the professor emphasizes:} real sustainability means changing the firm's \textbf{DNA} and technology (lecture example: Spanish fashion firm sourcing recycled plastics from a large network of fishermen).\\
\textbf{Strategic framing:} CSR vs substantive capability and supply-chain transformation.
}

\cornell{30) Vocabulary Word Bank (Risk, M\&A, Information) + CEO Ego Risk}
{Which terms show up in governance essays?\\
Why do many M\&As fail?}
{
\textbf{Beta (\(\beta\)):} stock risk vs the whole market (see block 26).\\
\textbf{M\&A (mergers \& acquisitions):} fast growth path; managers may favor it for status even when deals are expensive/risky for owners.\\
\textbf{Stock options:} contract to buy stock at a set price; aligns manager payoff with share-price performance if designed well.\\
\textbf{Agency problem:} inherent conflict when one party must act in another's interest (managers vs shareholders).\\
\textbf{Asymmetric information:} managers know more than shareholders; shifts bargaining power and enables hidden motives.\\
\textbf{Greenwashing:} deceptive environmental branding (see block 29).\\
\textbf{Strategic insight---CEO ego:} mergers often fail when CEOs deploy ''other people's money'': premiums of \textbf{20--50\%} over market value to build an empire---shareholders need strong \textbf{governance} (board) to limit managerial discretion.
}

\cornell{31) Shareholder Value + The Strategic Analysis Loop}
{How is success defined for owners?\\
What three questions structure strategy work?}
{
\textbf{Shareholder satisfaction:} value when \textbf{ROI} (or risk-adjusted return) exceeds \textbf{opportunity cost} of the next-best use of capital.\\
\textbf{Strategic analysis loop:} (1) \textbf{What} do we want?---mission/vision; (2) \textbf{Where} are we?---external and internal analysis; (3) \textbf{How} do we get there?---strategy formulation and implementation.\\
\textbf{Internal analysis role:} ensure the firm creates rather than destroys value relative to rivals and benchmarks.
}

\cornell{32) Porter Value Chain: Primary vs Support + Margin}
{How do you disaggregate the firm for analysis?\\
Where does margin come from?}
{
\textbf{Purpose:} split the firm into strategically relevant activities to trace \textbf{costs} and \textbf{sources of differentiation}.\\
\textbf{Primary activities:} inbound logistics, operations, outbound logistics, marketing/sales, service---direct creation, sale, and transfer of the product.\\
\textbf{Support activities:} firm infrastructure, HR management, technology development, procurement---enable primary activities.\\
\textbf{Margin:} difference between total value the buyer perceives and the collective cost of performing the activities (value created minus cost to create it).\\
\textbf{JIT (lean):} e.g.\ Toyota-style coordination with suppliers; produce/ship in tight sync with demand to cut inventory (order-driven flow).
}

\cornell{33) Linkages: Horizontal vs Vertical (Value System)}
{What ties activities together?\\
How do suppliers and channels fit?}
{
\textbf{Horizontal linkages:} connections \textit{inside} the firm (e.g.\ HR training improves operations quality/speed).\\
\textbf{Vertical linkages:} firm's chain linked to \textbf{suppliers'} and \textbf{distributors'/buyers'} chains---\textbf{value system} view.\\
\textbf{Exam move:} name one horizontal and one vertical linkage with a causal ''so what'' for cost or differentiation.\\
\textbf{Strategic groups:} clusters of rivals with similar strategies within an industry (e.g.\ luxury vs economy automakers); shapes competitive reference set.
}

\cornell{34) RBV + SWOT Internal Side (Strengths \& Weaknesses)}
{What is internal analysis for in SWOT?\\
What strategic rule ties S/W to action?}
{
\textbf{RBV lens:} strengths/weaknesses reflect resources and capabilities (heterogeneity, not just generic labels).\\
\textbf{Strengths:} core competencies that can underpin advantage when valuable, rare, and well deployed.\\
\textbf{Weaknesses:} performance below industry average or rivals in activities that matter for buyers.\\
\textbf{Strategic rule:} anchor strategy on \textbf{strengths}; use strategy to \textbf{neutralize, improve, or outsource} weaknesses rather than pretending they do not exist.
}

\cornell{35) Competitive Advantage: Efficiency vs Differentiation}
{What are the two classic ways to outperform?\\
What mechanisms back each?}
{
\textbf{Low cost / efficiency:} perform value-creating activities more cheaply than rivals---often \textbf{economies of scale}, \textbf{learning curves}, process discipline.\\
\textbf{Differentiation:} deliver unique perceived value so price can exceed rivals'---often \textbf{innovation}, quality, \textbf{brand}, service.\\
\textbf{Integration:} map each claim to the value chain activity where the mechanism lives (exam evidence).
}

\cornell{36) Vocabulary Word Bank (Value Chain \& RBV Support)}
{Quick definitions for short answers?}
{
\textbf{Opportunity cost:} return foregone on the best alternative use of the same capital (ties to shareholder value).\\
\textbf{Beta (\(\beta\)):} stock volatility vs market (see block 26).\\
\textbf{Outsourcing:} shifting non-core or weak activities to external specialists.\\
\textbf{Economies of scale:} lower unit cost as volume rises (fixed cost spread, purchasing power, etc.).\\
\textbf{Knowledge spillovers:} ideas and knowledge flow across people and firms---especially in \textbf{clusters} (e.g.\ Silicon Valley).\\
\textbf{Stock options:} right to buy shares at a set price; governance tool to align managers with owners (see block 27).
}

\cornell{37) Strategic Insight: Minimum Efficient Scale (e.g.\ ``500,000 Car Rule'')}
{Why do some industries have few large players?\\
What creates a barrier to entry?}
{
\textbf{MES (minimum efficient scale):} output level where average cost is low enough to compete; below it, unit costs stay punishing.\\
\textbf{Lecture example (auto):} order of magnitude \textbf{500,000 cars/year} at global scale to absorb heavy fixed costs (R\&D, plants, platforms); smaller volume often loses money per unit.\\
\textbf{Strategic implication:} technology and scale economics set ``rules of the game''---high \textbf{barriers to entry}; new entrants need scale, capital, or a radically different model.\\
\textbf{Exam caution:} treat the 500k figure as illustrative of MES logic, not a universal law across every segment.
}

\cornell{38) External Analysis: Sun Tzu, Four Pillars, and PESTEL}
{How does ``know enemy / know yourself'' map to tools?\\
What does PESTEL scan?}
{
\textbf{Strategic posture:} success blends external awareness with internal appraisal (Sun Tzu: know the enemy and yourself).\\
\textbf{Four pillars (Grant-style):} (1) simple consistent long-term goals; (2) understand competitive environment; (3) objective resource appraisal; (4) effective implementation (see also block 14).\\
\textbf{PESTEL:} macro forces---\textbf{P}olitical, \textbf{E}conomic, \textbf{S}ocio-cultural, \textbf{T}echnological, \textbf{E}nvironmental, \textbf{L}egal.\\
\textbf{Climate change:} often cuts across \textbf{E}nvironmental, \textbf{E}conomic, and \textbf{P}olitical/geopolitical dimensions---not one silo.
}

\cornell{39) PESTEL in Practice: Climate Change (Opportunities \& Threats)}
{How does a macro shift become strategy-relevant?\\
What are concrete illustration cases?}
{
\textbf{Opportunity thread:} environmental/technological change (e.g.\ Arctic ice melt) can open \textbf{new shipping routes} or access to resources (e.g.\ rare minerals framing around \textbf{Greenland} in lecture examples).\\
\textbf{Threat thread:} rising seas/coastal risk increases losses and \textbf{insurance} payouts; real estate and infrastructure face higher risk premia.\\
\textbf{Strategic action:} insurers (e.g.\ \textbf{Allianz}-style positioning) and \textbf{engineering/flood-defense} firms must adapt products, pricing, and exposure.\\
\textbf{Exam move:} label each force (P/E/S/T/E/L), then state \textit{mechanism} $\rightarrow$ \textit{O or T} $\rightarrow$ \textit{strategic response}.
}

\cornell{40) Industry Attractiveness, Barriers, and Substitutes}
{When is an industry ``attractive''?\\
What blocks entry? What counts as a substitute?}
{
\textbf{Attractiveness:} externally, when \textbf{opportunities exceed threats} for your firm; structurally, whether the industry allows \textbf{sustained high profitability} (Five Forces lens).\\
\textbf{Barriers to entry:} \textbf{economies of scale} (high volume lowers unit cost); \textbf{capital requirements} (fixed investment); \textbf{artificial barriers} (tariffs, licensing, political protection).\\
\textbf{Substitutes:} \textit{different technologies} that satisfy the \textbf{same customer need} (e.g.\ caffeine/energy products vs coffee for ``energy/alertness'').
}

\cornell{41) Porter Five Forces: Structure and Auto / EV Illustration}
{How does the model map to ``imperfect'' profit potential?\\
What does a mature industry example show?}
{
\textbf{Purpose:} diagnose competitive pressure and realistic earnings power (not assuming perfectly competitive markets).\\
\textbf{1 Rivalry:} often high in mature industries \textit{like global autos}---share fights via price, branding, capacity.\\
\textbf{2 New entrants:} often \textbf{low} when \textbf{MES} is huge (order-of-magnitude \textasciitilde 500k units/year lecture benchmark spreads fixed cost).\\
\textbf{3 Substitutes:} can rise with disruption---\textbf{EVs} vs internal combustion as technology/need-shift threat to incumbents' stacks.\\
\textbf{4 Buyer power:} high when many comparable brands and \textbf{low switching costs}.\\
\textbf{5 Supplier power:} high for \textbf{specialized} inputs---e.g.\ battery supply leverage in EV transition.
}

\cornell{42) Geographic Clusters: Silicon Valley, Boston, and Cluster Logic}
{Why does place still matter in a global economy?\\
What advantages accrue inside a cluster?}
{
\textbf{Definition:} geographic concentration of related firms, suppliers, and institutions---competition \textit{and} collaboration.\\
\textbf{Anchors:} world-class \textbf{universities} often seed talent and R\&D---\textbf{Stanford} + \textbf{UC Berkeley} (Silicon Valley); \textbf{MIT} + \textbf{Harvard} (Boston area).\\
\textbf{Benefits:} \textbf{labor pooling} (engineers switch employers with minimal relocation); \textbf{knowledge spillovers} (informal and formal diffusion of ideas); \textbf{specialized infrastructure} (legal, finance, suppliers tuned to the industry).\\
\textbf{Strategy link:} location choice and partnerships inside a cluster can raise innovation speed and access to complementary assets.
}

\cornell{43) External Vocabulary + ``Rules of the Game'' (Lobbying \& Regions)}
{Definitions for drills?\\
Why might identical ``business logic'' fail across countries?}
{
\textbf{Economies of scale:} unit-cost falls as volume rises (see blocks 37, 40).\\
\textbf{Greenwashing:} deceptive green branding (see blocks 11, 29, 30).\\
\textbf{Oligopoly:} few large sellers dominate (global aircraft or handset-style structures).\\
\textbf{Externalities:} costs/benefits borne by others without full price reflection (e.g.\ pollution).\\
\textbf{SWOT:} internal \textbf{S/W} vs external \textbf{O/T} synthesis (see block 34).\\
\textbf{Rules of the game:} competitive rivalry unfolds under \textbf{different institutional rules} by region (e.g.\ USA vs China); managers may need \textbf{lobbying} and policy engagement (Washington, Brussels) to shape \textbf{artificial barriers} that protect incumbents or segments---not only ``pure market'' tactics.
}

\cornell{44) Workshop Case: Madonna \& Internal Analysis (RBV Lens)}
{How does RBV apply to a person-as-brand?\\
What makes the case exam-useful?}
{
\textbf{Workshop focus:} collaborative internal analysis using the \textbf{Madonna} case to practice \textbf{Resource-Based View (RBV)} and \textbf{competitive advantage}.\\
\textbf{RBV takeaway:} strategy rooted in \textbf{resources and capabilities}---here, difficult-to-copy bundles (performance persona, vocal/artistic identity, reinvention rhythm).\\
\textbf{Non-imitable bundle:} combine talent signals (style, stagecraft) with \textbf{artistic control} and brand narrative so rivals cannot copy the \textit{system}, not only the song.\\
\textbf{Class aside (Gypsy roots):} brief historical/cultural frame (migration from India through Egypt/North Africa to Spain) used to contextualize certain musical ``soul''/roots traditions---link only if the prompt asks for cultural context.
}

\cornell{45) Grant Four Elements Applied (Goals, Environment, Resources, Implementation)}
{How do students map Grant to Madonna?\\
What evidence line should you write in an answer?}
{
\textbf{1 Goals:} simple, consistent, long-term (e.g.\ global stardom, not a two-year trend) so tactics align across decades.\\
\textbf{2 Environment:} read the music industry shift (radio $\rightarrow$ \textbf{MTV/visual media}) and exploit symbolic presence (high-production videos / visual branding).\\
\textbf{3 Resources (honest appraisal):} \textbf{tangible}---accumulated financial strength over time; \textbf{intangible}---\textbf{brand}, provocative positioning, accumulated reputation, governance of creative ``product.''\\
\textbf{4 Implementation:} repeated \textbf{reinvention} of image/sound to track demographic and cultural change while keeping a coherent star narrative.\\
\textbf{Exam structure:} one sentence per pillar tied to a \textit{case fact} $\rightarrow$ \textit{mechanism} $\rightarrow$ \textit{performance outcome}.
}

\cornell{46) Sustainable Advantage, Control, and ``CEO of Herself''}
{Why did advantage last when peers faded?\\
How does governance map to a solo star?}
{
\textbf{Sustainability:} not one hit---ongoing relevance via reinvention + control over artistic direction (core competency as \textbf{staying ahead of trends}).\\
\textbf{Agency \& control:} unusually strong say over artistic and business decisions---reduces dependence on label-only priorities and protects long-term brand logic.\\
\textbf{CEO metaphor:} treat the star as a \textbf{CEO} managing a portfolio of assets (music, visuals, tours, partnerships); success is \textbf{calculated strategy}, not only taste.\\
\textbf{Agency problem angle:} minimize misalignment by acting as her own strategic ``board''---controlling assets and direction so external agents do not capture value at the expense of long-run positioning.\\
\textbf{Peer contrast:} same industry shocks destroyed artists without comparable control/reinvention systems.
}

\cornell{47) Vocabulary Word Bank (Madonna / RBV Workshop)}
{Quick definitions from the exercise?}
{
\textbf{Inimitable:} hard for rivals to copy---key RBV language for durable differentiation.\\
\textbf{Sustainable development (workshop sense):} sustaining career and \textbf{brand value} over the long run (not the UN SDG list unless the prompt says so).\\
\textbf{Artistic direction:} top-level control over the creative ``product''---analogous to \textbf{R\&D + product definition} in a firm.\\
\textbf{Modest background:} starting point that can sharpen motivation and values (links to governance/motivation themes in earlier lectures).\\
\textbf{Implementation:} turning the strategic plan into repeated, observable actions and releases that hit targets.
}

\cornell{48) RBV Core: Resources, Capabilities, and Strategic Rarity}
{Why is RBV the ``engine'' beside PESTEL?\\
What categories of internal assets matter?}
{
\textbf{Core premise:} sustainable success depends on the \textbf{specific bundle of resources} the firm controls and how management \textbf{identifies and deploys} them (PESTEL sets context; RBV explains heterogeneity in performance).\\
\textbf{Tangible resources:} physical and financial---plant, equipment, inventory, cash; often visible on the balance sheet but \textbf{easier to buy or copy}.\\
\textbf{Intangible resources:} IP, \textbf{brand reputation}, \textbf{organizational culture}---often harder to replicate and therefore more strategic.\\
\textbf{Human resources / human capital:} skills, knowledge, motivation; includes \textbf{tacit knowledge} that is not fully codified.\\
\textbf{Organizational capabilities:} capacity to combine resources into routines (e.g.\ lean system, R\&D process, a ``hit-making'' collaboration like a band).\\
\textbf{Strategic rarity:} only \textbf{scarce} resources can underpin advantage; widely available inputs are ``table stakes,'' not a differentiator.
}

\cornell{49) VRIO: Testing for Sustainable Competitive Advantage}
{How do you know a resource is strategic?\\
What four filters does the professor use?}
{
\textbf{V --- Value:} does the resource exploit an \textbf{opportunity} or neutralize a \textbf{threat}?\\
\textbf{R --- Rarity:} is it controlled by only a few firms?\\
\textbf{I --- Imitability:} is it costly or complex to copy (history-dependent brand, tacit routines, causal ambiguity)?\\
\textbf{O --- Organization:} is the firm structured to \textbf{capture} value (incentives, coordination, complementary assets)?\\
\textbf{Exam use:} walk the four letters with one case asset; conclude: competitive parity vs temporary advantage vs sustained advantage (see also block 34 VRIO line).
}

\cornell{50) Resource Dynamics: Depreciation vs Appreciation}
{How do assets ``age'' differently?\\
What examples anchor the lecture?}
{
\textbf{Physical/financial:} tend to \textbf{depreciate}---machinery wears out; cash is consumed; maintenance and replacement matter.\\
\textbf{Intangible (reputation):} can \textbf{appreciate} with consistent delivery---trust and brand equity compound when reinforced.\\
\textbf{Human (experience):} often appreciates via the \textbf{learning curve}---speed and quality improve with deliberate practice (until obsolescence or turnover interrupts).\\
\textbf{Strategy implication:} invest in what compounds; manage obsolescence and retention for assets that walk out the door.
}

\cornell{51) Case Logic: Elite Soccer and ``Buying Stars'' vs Capability}
{What does football illustrate about rarity and teams?}
{
\textbf{Strategic rarity:} all clubs have players; few have world-class ``champions'' in the right roles at the right time.\\
\textbf{Capability beats shopping list:} winning needs \textbf{organization}---tactics, chemistry, leadership---not only individual talent.\\
\textbf{Purchase limit:} rivals can try to buy stars yet still underperform without compatible \textbf{culture}, systems, and integration (RBV: social complexity + tacit coordination).\\
\textbf{Transfer to firms:} M\&A or star-hiring without integration capability often destroys value.
}

\cornell{52) RBV Vocabulary + ``The Human Secret'' (Retention \& Collective Intelligence)}
{What terms explain why copying fails?\\
Why is human capital strategically awkward?}
{
\textbf{Tacit knowledge:} hard to capture in manuals or contracts; travels with people.\\
\textbf{Causal ambiguity:} rivals cannot see \textit{which} internal factors drive success, so imitation fails.\\
\textbf{Social complexity:} advantage rooted in trust, culture, reputation---not replicable by a single purchased input.\\
\textbf{Lean production:} systemic waste elimination (Toyota Production System / lean operations).\\
\textbf{Intellectual property (IP):} patents, copyrights, trademarks---legal protection for intangible creation.\\
\textbf{Human secret:} employees are not owned; they can exit with tacit know-how---winning firms build conditions where talent \textbf{chooses to stay} and feeds \textbf{collective intelligence}.
}

\cornell{53) Competitive Advantage: Cost Leadership vs Differentiation}
{How does a firm outperform rivals?\\
What is a price premium?}
{
\textbf{Competitive advantage:} systematically beat rivals on economics or buyer preference---not one-off luck.\\
\textbf{Cost leadership:} lowest-cost position in the strategic space scale, tight cost control, often ``no-frills'' offer (e.g.\ budget airline logic).\\
\textbf{Differentiation:} unique features, quality, or service buyers value; if \textbf{perceived value} is high, they accept a \textbf{price premium} that covers extra cost to differentiate.\\
\textbf{Link:} map each claim to activities in the value chain (block 32--35) and to evidence in the case.
}

\cornell{54) Strategic Clock: Price vs Perceived Value (Bowman)}
{What positions does the clock label?\\
Where do ``failure'' strategies sit?}
{
\textbf{Purpose:} position the firm on \textbf{price} and \textbf{perceived value}---extends beyond binary cost vs differentiation (see also block 5).\\
\textbf{Positions 1--2 (low price):} cost focus, volume, efficiency.\\
\textbf{Position 3 (hybrid):} high perceived value at relatively low price (e.g.\ IKEA-style ambition)---hard to sustain, potent if credible.\\
\textbf{Position 4 (differentiation):} strong value at standard or moderately higher price.\\
\textbf{Position 5 (focused differentiation):} luxury/focused premium---very high price matched by very high perceived value.\\
\textbf{Positions 6--8 (failure zone):} high price with \textbf{low} perceived value---unsustainable in competition; may persist only under market power.\\
\textbf{Exam tip:} state position, then why the firm is/not credible vs rivals.
}

\cornell{55) Market Segmentation \& Strategic Groups}
{Why tailor strategy to segments?\\
How do peer clusters matter?}
{
\textbf{Market segmentation:} split demand into groups with different needs (luxury vs economy, pro vs consumer, region, usage).\\
\textbf{Strategic implication:} choose segment, offer, and clock position that fit---avoid ``average'' products that satisfy nobody.\\
\textbf{Strategic groups:} rivals with similar models (fast fashion vs haute couture); defines direct comparators for price--value benchmarks.
}

\cornell{56) Industry Life Cycle: Stages \& Strategic Choices}
{How does industry maturity change the game?\\
What decisions dominate each stage?}
{
\textbf{Emergence / introduction:} high risk, low/zero profit, heavy spend to educate demand; core choice \textbf{entry timing} (lead vs follow).\\
\textbf{Growth:} market ``pie'' expands; room for many winners; often less zero-sum than later stages.\\
\textbf{Maturity:} growth slows; share is often \textbf{zero-sum}---to grow, take from rivals; consolidation via \textbf{M\&A} common.\\
\textbf{Decline:} shrinking demand; choices include \textbf{harvest} (milk cash flow), \textbf{niche} (loyal residual segment), or \textbf{exit}.\\
\textbf{Internationalization link:} when home market is mature/declining, firms may seek foreign growth---fits corporate/scope logic (block 5).
}

\cornell{57) Competitive Strategy Vocabulary + ``Monopoly Trap'' (Telefónica Example)}
{Key terms for short answers?\\
When can high-price/low-value survive?}
{
\textbf{Price premium:} extra willingness to pay for brand or differentiated benefits vs generic alternatives.\\
\textbf{Perceived value:} buyer's judgment of benefits relative to price (not the firm's internal cost story alone).\\
\textbf{Economies of scale:} lower unit cost as volume rises (reinforces low-price positions).\\
\textbf{M\&A:} consolidate to cut rivalry, add scale, or acquire capabilities/markets (maturity-stage tool).\\
\textbf{Internationalization:} expand abroad when domestic opportunity saturates or declines.\\
\textbf{Strategic groups:} see block 55.\\
\textbf{Monopoly trap:} high price + mediocre perceived value can persist only when customers have \textbf{no choice}. Lecture cites pre-liberalization \textbf{Telefónica} (Spain): limited competition allowed weak value--price bundles; entry of foreign rivals forced shift toward \textbf{cost} or \textbf{differentiation} credibility---``DNA'' change, not only pricing tweaks.
}

\cornell{58) Unit 6: Corporate Scope \& Growth Directions Matrix}
{What does ``where to compete'' add to ``how to compete''?\\
Which growth paths did the lecture contrast?}
{
\textbf{Corporate scope:} portfolio of businesses and activities the firm should own---\textbf{vertical integration} (upstream/downstream control) vs \textbf{diversification} (new products/markets).\\
\textbf{Beyond business-level rivalry:} corporate strategy chooses arenas; business strategy wins inside an arena (see block 5 integration sentence).\\
\textbf{Growth directions (lecture matrix):}\\
\textbf{Specialization / same business:} deepen share in one market---lower strategic risk; often powered by \textbf{economies of scale}.\\
\textbf{Diversification:} new products and/or markets---higher risk; needs new capabilities or M\&A.\\
\textbf{Vertical integration:} buy suppliers or distributors---more control, less flexibility if mis-sized.\\
\textbf{Internationalization:} new countries---growth option when home demand saturates; legal/cultural friction.
}

\cornell{59) Case Lens: Starbucks, Third Place, Clustering, and Differentiation}
{What bundle justifies a price premium?\\
What is ``third place'' in strategy terms?}
{
\textbf{Third place:} social setting between home (first) and work (second)---\textbf{atmosphere, community, routine}---core to high \textbf{perceived value} in the Starbucks case framing.\\
\textbf{Starbucks logic:} commodity-like coffee augmented by \textbf{intangibles}---brand, convenient locations, store ambiance, partner/employee experience cues---supporting \textbf{differentiation} and premium pricing (lecture cited an illustrative on-the-order-of-\textbf{six-dollar} cup; verify to your slide).\\
\textbf{Clustering:} dense store placement in cities to build habit, visibility, and operational learning---corporate expansion tactic tied to scale and brand presence.\\
\textbf{Exam line:} separate \textbf{physical product} similarity from \textbf{buyer value} created by the full experience.
}

\cornell{60) Scale, Experience Curve, Diversification Types, and M\&A Premium}
{What economic mechanisms support staying ``specialized''?\\
When is diversification related vs unrelated?}
{
\textbf{Economies of scale:} higher volume spreads fixed costs---unit cost falls; supports leadership in one core category.\\
\textbf{Experience curve:} repetition cuts errors and cycle time---learning-driven efficiency over time.\\
\textbf{Diversification---related:} adjacencies linked to core competences (lecture: Starbucks-style \textbf{machines/grocery channels} extending reach without leaving the brand logic).\\
\textbf{Diversification---unrelated:} moves far from core---often \textbf{value-destructive} without integration skills.\\
\textbf{M\&A for growth:} fastest but costly; buyers may pay a \textbf{control premium} on the order of \textbf{30--50\%} above market price to win control (treat as lecture band, not a universal law).
}

\cornell{61) Industry Life Cycle: Starbucks North America vs International Growth}
{Why expand to China and India?\\
How does life-cycle stage drive corporate moves?}
{
\textbf{North American coffee retail:} portrayed as approaching \textbf{maturity/decline}---saturated shops, slower organic growth.\\
\textbf{Strategic necessity:} to keep corporate growth, pursue markets where the \textbf{same industry} is earlier in the cycle (\textbf{growth} in China/India story in lecture).\\
\textbf{Tie-in:} links block 56 (life cycle) to block 58 (internationalization)---\textbf{scope} follows where incremental demand lives.\\
\textbf{Case narrative:} maturity at home pushes \textbf{geographic scope}; execution risk includes local competition, regulation, and brand translation.
}

\cornell{62) Stuck in the Middle, Porter Generics, and the Zara Hybrid Exception}
{What does ``stuck in the middle'' mean?\\
When can a hybrid still work?}
{
\textbf{Stuck in the middle (Porter):} no credible \textbf{cost leadership} and no sharp \textbf{differentiation}---muddled trade-offs; often underperforms firms with clear positioning (see blocks 53--54).\\
\textbf{Zara / Inditex exception:} \textbf{hybrid}---tight cost/operations (scale, logistics, fast replenishment) plus fashion-speed and trend-led assortment (differentiation on responsiveness); requires rare operating model, not generic ``average'' pricing.\\
\textbf{Bowman clock (lecture recap):} positions \textbf{1--2} low price (cost/no-frills); \textbf{4--5} high perceived value (differentiation / focused differentiation); \textbf{6--8} failure band (high price, low perceived value) except where temporary market power or convenience traps apply.\\
\textbf{Exam discipline:} if claiming hybrid, show \textit{two} coherent mechanisms (cost driver + differentiation driver) backed by case facts.
}

\cornell{63) Business Development: Diversification \& Vertical Integration}
{Related vs unrelated diversification?\\
Backward vs forward integration?}
{
\textbf{Related diversification:} adjacent products/markets that reuse competences or assets (lecture: car firm $\rightarrow$ \textbf{car insurance}).\\
\textbf{Unrelated diversification:} ``conglomerate'' moves far from core---higher integration and governance risk; easier to destroy value.\\
\textbf{Vertical integration:} own more steps in the value chain---\textbf{backward} (supplier/producer upstream) vs \textbf{forward} (distributor, retail, direct channels downstream).\\
\textbf{Trade-off:} coordination and margin capture vs capital intensity and loss of supplier flexibility (block 58).
}

\cornell{64) Strategic Clock: Supermarkets (Mercadona) and Convenience ``Failure'' Positions}
{How does grocery retail map to price vs perceived value?\\
When is a high price ``excused''?}
{
\textbf{Low price (1--2):} discount format---basic assortment, minimal service frills.\\
\textbf{Hybrid (3):} strong perceived value at disciplined price---lecture cites \textbf{Mercadona} (Spain) as exemplar.\\
\textbf{Differentiation (4):} broader service/assortment at higher price.\\
\textbf{Focused differentiation (5):} gourmet/premium food retail.\\
\textbf{Positions 6--8:} weakness when price is high and perceived value is low---lecture: some \textbf{small convenience} shops sit toward \textbf{position 7} (failure zone) because they charge convenience prices for modest perceived value yet win traffic from \textbf{24/7} or late-night availability (e.g.\ when larger stores are closed).\\
\textbf{Strategic insight---price of convenience:} that local ``only option at 2:00 a.m.'' logic is \textbf{fragile}: scaled retailers with logistics can undercut the trap; durable strategy still requires moving up the clock via lower price or materially stronger customer experience.
}

\cornell{65) Airline Positioning: LCC vs Full-Service Differentiation}
{How do budget and network carriers map to the clock?}
{
\textbf{Low-cost carriers (LCC):} aim near clock \textbf{1--2}---no frills, dense seating, secondary airports where applicable, relentless unit-cost discipline.\\
\textbf{Full-service carriers:} differentiation (around position \textbf{4})---lounges, loyalty/Frequent-flyer programs, service quality cues to earn a \textbf{price premium}.\\
\textbf{Exam move:} tie each bundle to buyer segment and switching costs; avoid labelling all airlines with one generic strategy.
}

\cornell{66) Ansoff Growth Paths (Penetration, Development, Diversification)}
{How does the matrix structure ``what to sell'' vs ``where''?\\
What examples anchor the lecture?}
{
\textbf{Market penetration:} more of the \textbf{same} product to the \textbf{same} market---push share, outlets, frequency (home-market sales growth).\\
\textbf{Market development:} \textbf{same} product, \textbf{new} geography---lecture: \textbf{Inditex/Zara} Spain $\rightarrow$ Mexico/China-style expansion.\\
\textbf{Product development:} \textbf{new} product, \textbf{same} customers---lecture: \textbf{Starbucks Via} instant coffee extension.\\
\textbf{Diversification:} \textbf{new} product and \textbf{new} market---lecture: construction firm $\rightarrow$ energy business (unrelated-style illustration).\\
\textbf{Link:} pairs with blocks 58 (corporate scope) and 60 (related vs unrelated logic).
}

\cornell{67) Positioning \& Growth Vocabulary (Hybrid, MES, Asset Stripping, Synergy)}
{Definitions for exam fragments?}
{
\textbf{Hybrid strategy:} simultaneous cost discipline and meaningful differentiation (lecture cites \textbf{Toyota} or \textbf{Zara}-type ambition---requires integrated operating choices).\\
\textbf{Perceived value:} buyer-side benefit vs price judgment (see blocks 53, 57, 64).\\
\textbf{Minimum efficient scale (MES):} smallest scale at which unit costs approach the industry frontier---barrier logic in blocks 37, 41.\\
\textbf{Asset stripping:} buy a target and sell parts for cash---often tied to \textbf{unrelated} or financial plays; governance/reputation risk.\\
\textbf{Synergy:} combined value exceeds the sum of standalone parts (``\textbf{1 + 1 = 3}'' shorthand in lecture).
}

\section*{Finals Tile Insights (From Uploaded Images)}
\begin{itemize}
  \item The Starbucks case tiles reinforce a high-value exam pattern: \textbf{strategy success -> crisis -> renewal -> forward risk}.
  \item First-midterm feedback indicates need for tighter Five Forces depth and less generic phrasing in recommendations.
  \item Strong answers in this course should explicitly connect framework terms to \textbf{case evidence} before concluding.
  \item Prioritize ``how to compete'' plus ``where to compete'' integration in final long answers.
\end{itemize}

\section*{Summary Block (Cornell Bottom)}
\begin{itemize}
  \item Strategic Management finals answers should be mechanism-led, evidence-based, and structured.
  \item Reuse this Cornell format as a transferable study template for high-risk courses.
  \item Keep closing line in every answer: strategic implication plus condition for success.
\end{itemize}

\end{document}
